% ============================================================
%   Security frontend rationale
% ============================================================

\documentclass[11pt,a4paper]{article}

\usepackage[a4paper, margin=22mm]{geometry}
\usepackage[T1]{fontenc}
\usepackage{lmodern}
\usepackage{microtype}
\usepackage{xcolor}
\usepackage{booktabs}
\usepackage{tabularx}
\usepackage{graphicx}
\usepackage{enumitem}
\usepackage{hyperref}
\hypersetup{
    colorlinks=true, linkcolor=accent, urlcolor=accent,
    citecolor=accent, breaklinks=true,
}

\definecolor{accent}{RGB}{20,70,140}

\setlength{\parskip}{4pt}
\setlength{\parindent}{0pt}

\newenvironment{question}[1]{%
    \par\smallskip
    \noindent\textcolor{accent}{\bfseries Q.\ #1}\par\nopagebreak
    \smallskip\noindent\ignorespaces
}{\par\medskip}

\title{\color{accent}\bfseries
       Why the EPSS-TPG security frontend looks the way it does\\[2pt]
       \large A walk-through of the rationale, an inventory of
       what the security frontend adds, and the
       questions we are answering or might face}
\author{}
\date{}

\begin{document}
\maketitle
\thispagestyle{empty}

% ============================================================
\section*{What this document is, and why it exists}

The EPSS-TPG pipeline turns the CVE input text (the NVD
description, an LLM-generated summary, or both joined
together) into a Text Property Graph and feeds that graph
to a multi-view GNN. Most of the graph comes straight from
spaCy: tokens, sentences, noun phrases, named entities,
and the syntactic edges between them. Once the spaCy layer
is in place, a second stage scans the same text with a
small library of regular expressions and curated
dictionaries and lays a layer of security nodes and edges
on top of the linguistic backbone. This second stage is
the \emph{security frontend}.

Three reasonable questions tend to come up the first time
someone reads about it:

\begin{enumerate}[leftmargin=*, itemsep=0pt]
    \item Why bother with rule-based extraction at all?
          A neural Named Entity Recognition (NER) model
          would surely be smarter.
    \item What exactly is added on top of the spaCy graph,
          and what is the basis for picking those particular
          entities and edges?
    \item Which keywords are hard-coded, and do they
          actually matter for the downstream model?
\end{enumerate}

This note answers all three, and then collects the
follow-up questions we tend to come back to whenever the
design is challenged or revisited. Every count and every
inventory in this document was re-checked against the
implementation on the day of writing.

% ============================================================
\section{Why a rule-based security layer is worth the trouble}
\label{sec:why-bother}

There are three answers, and they reinforce one another.

\paragraph{1.\ The most useful labels are easy to write as
rules.}
The single most useful thing the security frontend extracts
is the CVE identifier itself. A CVE-ID has exactly one shape:
the literal letters CVE followed by a year and a sequence
number. A regex matches it perfectly on every NVD advisory
ever written. The same is true of CWE identifiers, CVSS
scores, and semantic version strings. These are
syntactically rigid. Training a model to learn them would
only add noise: a neural NER might miss one, mislabel one,
or hallucinate one that is not there. A regex never does.
Spending a few milliseconds per CVE on a regex match is a
much better deal than spending a tenth of a parameter budget
on relearning what the specification already pins down.

The same logic extends, with somewhat softer precision, to
the curated dictionaries. Software product names (OpenSSL,
nginx, Log4j) are a closed and slow-moving vocabulary.
Common attack-vector phrases (\emph{remote},
\emph{unauthenticated}, \emph{user-supplied input}) and
impact phrases (\emph{remote code execution}, \emph{denial
of service}, \emph{privilege escalation}) recur across
thousands of advisories with nearly identical surface forms.
A hand-maintained dictionary captures them with high
precision and is small enough that an analyst can audit it
in an afternoon.

\paragraph{2.\ A rule-produced graph can be audited
end-to-end.}
Every entity created by the security frontend carries a
tag identifying which dictionary or regex produced it,
together with the character offsets in the original text.
When the GNN flags a CVE as high-risk and a human asks
``which evidence in the text drove that?'', the answer is
reachable: it was these specific entity nodes, each
created by this specific match. A learned extractor would
make the same decisions, but the audit trail would
disappear. That trade-off only becomes attractive once
recall becomes a real bottleneck, and a separate hybrid
layer based on contextual-embedding similarity addresses
recall without giving up the audit trail.

\paragraph{3.\ Typed security edges are what give the
encoder a fifth view.}
The multi-view GGNN encodes the graph one edge family at a
time and then fuses the views by attention. Without typed
security edges, every security relation would collapse into
a generic ``entity-to-entity'' bucket and the encoder would
have no way to isolate a security view at all. The edge
typing is not just bookkeeping; it is what makes the fifth
view of the encoder possible.

% ============================================================
\section{Why spaCy or SecBERT alone are not enough}
\label{sec:why-not-spacy}

This is the question that comes up most often.

spaCy's English NER produces tags from the eighteen-label
OntoNotes~5 scheme: \textsc{cardinal, date, event, fac,
gpe, language, law, loc, money, norp, ordinal, org,
percent, person, product, quantity, time, work\_of\_art}.
None of those mean anything to a security model.
Concretely:

\begin{itemize}[leftmargin=*, itemsep=0pt]
    \item \texttt{CVE-2024-1234} becomes a plain
          \textsc{token}. spaCy has no schema slot for it.
    \item \emph{OpenSSL}, \emph{nginx} and \emph{Log4j} may
          or may not be tagged \textsc{org} or
          \textsc{product} depending on the surrounding
          context. Even when they are, there is no
          \textsc{software} label that the downstream GNN
          can attend to as its own slot.
    \item \emph{CVSS 9.8}, \emph{critical severity},
          \emph{buffer overflow}, \emph{RCE} carry no
          special meaning to a generic parser.
    \item spaCy never emits security \emph{relations}
          (CVE\,$\to$\,CWE, CVE\,$\to$\,Software,
          Attack-Vector\,$\to$\,Impact). The whole typed
          security edge layer would be missing.
\end{itemize}

SecBERT is closer to the right tool, but it is a contextual
embedding model, not a tagger. The pipeline uses it
extensively: it embeds every token, sentence, and entity
node with a 768-dimensional contextual vector, and that is
where most of the model's semantic capacity comes from.
What it does not do on its own is decide that ``this span is
a CVE-ID, that span is a software product, that other span
is an attack vector''. That classification step is still
the job of either a tagger (neural or rule-based) or a
downstream model trained on labelled spans. Building a
labelled-span dataset for CVE descriptions is its own
multi-month project. Rules give roughly ninety percent of
the signal for almost no annotation cost, and the hybrid
fallback adds a contextual-embedding step for the rest.

So the layered design, in plain English, is this: spaCy
provides the syntactic backbone, SecBERT provides the
contextual semantics on every node, the security frontend
provides the security-specific labels and edges, and the
hybrid layer recovers the entities the rules missed.

% ============================================================
\section{What the security frontend actually adds to the graph}
\label{sec:inventory}

\subsection{Twelve security node types in the schema,
            ten of them populated}

The schema declares twelve security node types:

\begin{center}\small
CVE\_ID, CWE\_ID, SOFTWARE, VERSION, COMPONENT,
CODE\_ELEMENT, ATTACK\_VECTOR, IMPACT, REMEDIATION,
THREAT\_ACTOR, VULN\_TYPE, SEVERITY
\end{center}

Ten of these are populated by the extractors: CVE\_ID,
CWE\_ID, SOFTWARE, VERSION, CODE\_ELEMENT, ATTACK\_VECTOR,
IMPACT, REMEDIATION, VULN\_TYPE, SEVERITY. The remaining
two (COMPONENT and THREAT\_ACTOR) are reserved slots in the
schema. They are not yet emitted, because the corresponding
vocabularies have not been curated. The honest count is
``ten populated security node types out of twelve declared''.

In the graph itself, every security entity is stored as an
ordinary entity node. The ``security-ness'' lives in two
properties: a top-level label (CVE\_ID, SOFTWARE, and so on)
and a narrower tag (\emph{rce}, \emph{remote},
\emph{buffer\_overflow}, and so on). The unified label
space prefixes the security labels with \texttt{SEC\_*} so
that the GNN sees a single integer slot per security type.

\subsection{Ten security edge types, all of them emitted}

The schema also declares ten typed security edges, and the
security-relations stage emits each one whenever both
endpoint entities exist in the same document
(\autoref{tab:sec-edges}). The basis column points at the
canonical security-data field each edge mirrors.

\begin{table}[!htbp]
\centering\small
\begin{tabularx}{\linewidth}{l l l X}
\toprule
\textbf{Edge} & \textbf{Source} & \textbf{Target} &
\textbf{Basis (canonical security data)} \\
\midrule
SEC\_AFFECTS         & CVE              & SOFTWARE       &
    NVD ships an explicit affected-products list per CVE. \\
SEC\_HAS\_VERSION    & SOFTWARE         & VERSION        &
    NVD ships affected version ranges per product. \\
SEC\_CLASSIFIED\_AS  & CVE              & CWE            &
    NVD assigns one or more CWE weakness classes to every
    CVE. \\
SEC\_LOCATED\_IN     & VULN\_TYPE       & CODE\_ELEMENT  &
    Where the weakness sits in code, for example a buffer
    overflow inside a string-copy primitive. \\
SEC\_EXPLOITED\_BY   & CVE              & ATTACK\_VECTOR &
    The attack-vector axis of CVSS (network, local,
    physical, adjacent). \\
SEC\_CAUSES          & ATTACK\_VECTOR   & IMPACT         &
    The impact axis of CVSS (confidentiality, integrity,
    availability) mapped to RCE, DoS, information
    disclosure, and so on. \\
SEC\_MITIGATED\_BY   & CVE              & REMEDIATION    &
    Mitigation, upgrade, or patch language at the end of
    most advisories. \\
SEC\_USES\_FUNCTION  & CVE              & CODE\_ELEMENT  &
    Direct mention of a vulnerable function or construct in
    the description. \\
SEC\_THREATENS       & ATTACK\_VECTOR   & SOFTWARE       &
    Network- or remote-reachable vectors threaten the
    product they can reach. \\
SEC\_HAS\_SEVERITY   & CVE              & SEVERITY       &
    CVSS's qualitative bands (critical, high, medium, low,
    informational). \\
\bottomrule
\end{tabularx}
\caption{The ten typed security edges and the canonical
security-data field each one mirrors. The full edge
vocabulary the GNN sees is therefore 13 base edge types
plus these 10 security edge types, for 23 in total when the
security view is active.}
\label{tab:sec-edges}
\end{table}

\subsection{Eight regex patterns and one regex list}

The security frontend defines eight specific regex patterns
plus the remediation pattern list
(\autoref{tab:regex}).

\begin{table}[!htbp]
\centering\small
\begin{tabularx}{\linewidth}{l X}
\toprule
\textbf{Pattern} & \textbf{What it matches} \\
\midrule
CVE identifier &
    \texttt{CVE-YYYY-NNNN+}, including the six Unicode
    hyphen variants in the range U+2010 to U+2015 that
    LLM-generated summaries often emit. \\
CWE identifier &
    \texttt{CWE-NNNN}. \\
Semantic version &
    Versions of the form \texttt{1.2.3}, \texttt{1.0.0-rc1},
    or \texttt{4.18.0-372}. A context filter rejects matches
    that look like CVSS scores. \\
CVSS score &
    \emph{CVSS 9.8}, \emph{CVSS:9.8}, \emph{CVSS v3 9.8},
    \emph{9.8/10}. \\
Severity with context &
    Adjective plus security-context noun (\emph{critical
    severity}, \emph{high CVSS score}). The mandatory
    second token blocks false positives like ``high
    availability''. \\
Bare-severity fallback &
    A bare severity word followed by punctuation, fired
    only when a security keyword is within sixty characters. \\
Code element &
    Sixteen dangerous C-standard-library and scripting
    function names (\emph{strcpy, strcat, sprintf, gets,
    scanf, memcpy, memmove, malloc, free, printf, fprintf,
    eval, exec, system, popen, fork}) plus the generic
    ``name()'' shape. \\
Vendor-product fallback &
    Capitalised \emph{Vendor Product} fragments immediately
    before a version number. Catches niche products absent
    from the curated software dictionary. \\
Remediation actions &
    A list of sixteen small patterns mapped to the tags
    \emph{upgrade, patch, update, workaround, mitigation,
    disable\_feature, restrict\_access, network\_block}. \\
\bottomrule
\end{tabularx}
\caption{Regex patterns in the security frontend. Each
pattern is paired, where appropriate, with a context filter
that rejects false positives in surrounding sentences.}
\label{tab:regex}
\end{table}

\subsection{Five keyword dictionaries}

\autoref{tab:dict-counts} lists the five keyword
dictionaries, their current sizes, and a sample of the
narrow tags each one emits. The counts come from a direct
tally of the literal entries in the implementation.

\begin{table}[!htbp]
\centering\small
\begin{tabularx}{\linewidth}{l r X}
\toprule
\textbf{Dictionary} & \textbf{Phrases} &
\textbf{Examples and narrow tags} \\
\midrule
Software product names & 127 &
    apache, nginx, openssl, openssh, linux kernel, mysql,
    postgresql, log4j, struts, jackson, fastjson, vmware,
    vcenter, esxi, fortios, pan-os, asa, big-ip, moveit,
    confluence, jira, exchange server, $\dots$ \\
Attack vectors & 47 &
    \emph{remote, network, local, physical, adjacent,
    user input, crafted input, malicious input, attacker} \\
Impact types & 51 &
    \emph{rce, dos, info disclosure, privesc, auth bypass,
    sqli, xss, csrf, ssrf, uaf, buffer overflow,
    integer overflow, heap overflow, stack overflow,
    path traversal, data breach, data exposure} \\
Vulnerability types & 47 &
    \emph{buffer overflow, buffer over-read, use after free,
    double free, race condition, command injection, sqli,
    xss, csrf, ssrf, xxe, deserialization, unsafe reflection,
    path traversal, missing authentication, weak cryptography,
    hardcoded credentials, prototype pollution, redos} \\
Code constructs & 44 &
    serialization, deserialization, reflection api, object
    graph, template engine, yaml/xml/json/url parser,
    request handler, session manager, session token,
    prepared statement, csrf token, authorisation check,
    $\dots$ \\
\bottomrule
\end{tabularx}
\caption{Sizes of the five keyword dictionaries used by the
security frontend.}
\label{tab:dict-counts}
\end{table}

% ============================================================
\section{Why these specific entities and edges, and not others}
\label{sec:why-these}

The inventory above is opinionated, and we are aware of
that. The following paragraphs explain each major design
choice, framed in terms of what the downstream
exploitation-prediction task actually cares about.

\paragraph{Entities are picked for what the model can act
on, not for what is easy to extract.}
VULN\_TYPE (what the bug is) and IMPACT (what the attacker
gets) are central because they predict exploitation
directly: memory-corruption bugs and RCE outcomes correlate
strongly with high EPSS scores; denial-of-service outcomes
correlate much less. ATTACK\_VECTOR captures reachability,
where network- and remote-without-authentication wins.
SOFTWARE and VERSION let the model attend to
product-specific patterns; a Log4j-class bug is not an
obscure-library bug. REMEDIATION is a coarse proxy for
``is a patch available'', which shifts the exploitation
timeline. SEVERITY and CVSS scores are the assessor's
qualitative band, complementing the numerical CVSS features
in the tabular branch. CVE\_ID and CWE\_ID anchor the
description to the NVD and CWE ontologies; even one CVE
identifier in the text tells the model it is reading a
vulnerability description.

\paragraph{Edges connect entities that CVE descriptions
actually relate.}
The ten security edges are not invented. Each one mirrors
a field that already exists in NVD or CVSS. ``CWE to
Mitigation'', for example, was deliberately not added as a
typed edge, because CVE advisories almost never write a
mitigation in terms of a CWE category; they write it in
terms of the specific CVE. The security frontend is
therefore deliberately conservative: ten edge types, all
with a clear data-model origin, all checkable against an
NVD record.

% ============================================================
\section{Questions we are answering, and ones we expect to face}
\label{sec:faq}

\begin{question}{How accurate is this? Has anyone actually
    measured the precision and recall?}
The high-precision regex extractors (CVE-ID, CWE-ID,
semantic version, CVSS) are designed to achieve effectively
one-hundred-percent precision on NVD-formatted text by
construction; the shapes they match are syntactic. The
dictionary extractors trade some precision for coverage, but
each one is anchored on word boundaries and uses
longest-first matching to avoid the classic ``git inside
github'' failure mode. A labelled benchmark of
human-annotated CVE spans is not yet available, so a precise
P/R number against ground truth cannot be quoted today. What
can be said is that the extractor is deterministic and
auditable, so any failure can be traced to a specific rule
and fixed in one place. Building a labelled test set is
the sensible next step the day such a benchmark is needed.
\end{question}

\begin{question}{Doesn't the version regex confuse CVSS
    scores like 9.8 with software versions?}
That is the failure mode the version extractor is
specifically guarded against. The version pattern is paired
with a context filter that rejects any candidate whose
surface shape looks like a CVSS score \emph{and} whose
surrounding sentence (within plus or minus eighty
characters) mentions \emph{CVSS}, \emph{base score} or
\emph{severity score}. A second short-window guard rejects
the ``9.8/10'' idiom. Severity uses a similar two-stage
design: a precise pattern that requires a security-context
noun runs first, and the bare-severity fallback only fires
when a security keyword sits within sixty characters. The
general principle is that every loose pattern is paired
with a context filter.
\end{question}

\begin{question}{Why not just train a neural NER on CVE
    descriptions and skip the rules entirely?}
Three reasons. First, there is no widely-used labelled
corpus of CVE spans, so the team would have to annotate one,
which is months of work. Second, the dominant entities
(CVE-IDs, CWE-IDs, versions, CVSS scores) are exactly the
cases where rules win by construction; a neural NER would
have to learn them and could hallucinate. Third, the
audit-trail story matters. Anyone asking ``why did this
CVE receive a high predicted score?'' can follow the
chain back to specific entity nodes and the specific regex
match that produced each one. With a neural NER, the chain
breaks at the model boundary. The combination of
high-precision rules plus a contextual-embedding fallback
gives most of the recall benefit without giving up the
audit trail.
\end{question}

\begin{question}{Where did the dictionary entries come
    from in the first place?}
The vocabulary grew empirically while running the pipeline
over the social-media dataset, the Megavul dataset, and
the NVD corpus, and inspecting which entity types
under-fired. The software list started
from the products that recur most often in NVD descriptions
and was extended whenever a recurring CVE family (Telerik
UI, FortiOS, vCenter, Confluence, MoveIt, Log4j, and the
like) was missed. The impact, vulnerability-type and
attack-vector dictionaries came from standard CWE and OWASP
terminology with duplicates merged. The remediation
patterns came from inspecting the closing sentences of
advisories. The vendor-product fallback regex catches niche
proprietary products that never reach the curated
dictionary.
\end{question}

\begin{question}{What happens when a phrase like ``buffer
    overflow'' fits two dictionaries at once?}
The three keyword extractors that compete for shared
phrases (VULN\_TYPE, IMPACT, ATTACK\_VECTOR) share a single
list of already-claimed character ranges. VULN\_TYPE runs
first, because it is the most specific tag. IMPACT runs
next, and ATTACK\_VECTOR last. The first extractor to claim
a character range wins. About fourteen phrases overlap
between IMPACT and VULN\_TYPE; they all end up tagged as
VULN\_TYPE.
\end{question}

\begin{question}{Is the security view actually used in the
    main training runs, or is this just on paper?}
The nineteen-run social-media ablation and the fifteen-run
Megavul ablation reported elsewhere in this work were
trained without the security-edge view, on the base
thirteen-edge graph. The security-edge view was added to
the pipeline later. The NVD/KEV temporal runs and the most
recent Megavul runs do use it. The honest framing for a
paper is therefore: the security view is an opt-in fifth
view; ablations are reported both ways. The figures that
say ``four views (or five with security)'' reflect this.
\end{question}

\begin{question}{How much upkeep does this need? When NVD
    adds a new CWE class, does anything have to change?}
Almost nothing. The CWE regex matches any
\emph{CWE-NNNN}-shaped string regardless of which CWE
number it carries, so new CWE classes appear in the graph
automatically. The software dictionary is the one piece
that ages: when a new vendor product starts showing up in
advisories and the vendor-product fallback misses it, an
analyst adds the name to the dictionary. This amounts to a
few minutes per quarter in practice.
\end{question}

\begin{question}{What about CVE descriptions written in
    languages other than English?}
The current extractors assume English text. The dictionaries
are in English and the spaCy model is the English
small-web model. CVE descriptions in NVD are English by
policy. Translated text from the LLM-summary pipeline is
also English-only. If a non-English corpus is ever
ingested, a language-detection step plus a per-language
frontend would be the right pattern, and the schema already
supports stacking frontends.
\end{question}

\begin{question}{If the rules already do most of the work,
    what is the hybrid layer for?}
The hybrid layer runs the rule-based extractor first, then
uses contextual-embedding similarity to find spans that are
close in embedding space to existing security entities but
that the rules missed. Rule entities win on overlap, so
the audit trail stays intact for the rule-detected
entities. This is the recall safety net: it picks up novel
software names and re-phrasings the dictionaries do not yet
know about.
\end{question}

\begin{question}{Does upper-case versus lower-case matter?
    Will the extractor miss ``OPENSSL'' written in all
    capitals?}
No. Every regex is case-insensitive and every dictionary
lookup is performed against the lowercased version of the
text. The original casing is preserved on the resulting
entity node so the display still reads naturally.
\end{question}

\begin{question}{Why does any of this belong in a paper?
    Isn't it just pre-processing?}
Two reasons. First, the multi-view encoder's security view
only exists because of these typed edges; turning it off
collapses the architecture from five views to four. The
extractor is therefore part of the model description, not
just a data loader. Second, the model is interpretable in
part because this pre-processing is auditable. A paper
that simply ``fed text to a big GNN and reported numbers''
would not need this section; this one is not that paper.
\end{question}

% ============================================================
\section{Empirical evidence: the GPT ablation (five variants)}
\label{sec:gpt-ablation}

To check whether the security frontend actually earns its
keep against the new 15-run social-media baseline
(Section~\ref{sec:smp-variant-comparison}), the model was
re-trained on the same five GPT text-source variants
(\emph{D}, \emph{SMP}, \emph{S\_git}, \emph{S\_cvss},
\emph{ALL}) with the security frontend disabled. In this
ablation the graph is built by the plain spaCy frontend
only, so every security entity node and every typed
security edge is removed; the NOSEC runs are kept in a
separate cache so they do not collide with the WITH-sec
runs. Everything else
(architecture, hyperparameters, data splits, soft EPSS
regression target, EPSS feature removed from the tabular
branch to prevent leakage) is identical. The test sets are
identical between the two halves of each comparison
(\autoref{tab:gpt-test-sizes}), so the numbers are directly
comparable.

\begin{table}[!htbp]
\centering\small
\begin{tabular}{lrrr}
\toprule
Variant & Test CVEs & Positives & Positive \% \\
\midrule
\emph{D}       & 1\,385 & 215 & 15.5\,\% \\
\emph{SMP}     & 1\,383 & 214 & 15.5\,\% \\
\emph{S\_git}  &   339 &  58 & 17.1\,\% \\
\emph{S\_cvss} & 1\,385 & 215 & 15.5\,\% \\
\emph{ALL}     & 1\,385 & 215 & 15.5\,\% \\
\bottomrule
\end{tabular}
\caption{Test-set composition for the five GPT variants
(soft EPSS target). The same test split is used for both
the WITH-security and NOSEC runs of each variant.
\emph{SMP} loses two CVEs to an empty social-media post
field; \emph{S\_git} evaluates on the subset of CVEs
with a non-empty GitHub-references summary.}
\label{tab:gpt-test-sizes}
\end{table}

\autoref{tab:gpt-headline} shows the headline ranking
metric (test PR-AUC) for each variant, side by side, and
\autoref{tab:gpt-full} expands the comparison to the full
six-metric table.

\begin{table}[!htbp]
\centering\small
\begin{tabular}{lrrrl}
\toprule
Variant & WITH security & NOSEC & $\Delta$ PR-AUC & Winner \\
\midrule
\emph{D}       & \textbf{0.8394} & 0.5416          & $-0.2978$ ($-35.5\,\%$) & WITH  \\
\emph{SMP}     & \textbf{0.8321} & 0.5738          & $-0.2583$ ($-31.0\,\%$) & WITH  \\
\emph{S\_git}  & 0.7553          & \textbf{0.7846} & $+0.0293$ ($+3.9\,\%$)  & NOSEC \\
\emph{S\_cvss} & 0.4916          & \textbf{0.5819} & $+0.0903$ ($+18.4\,\%$) & NOSEC \\
\emph{ALL}     & \textbf{0.7878} & 0.5286          & $-0.2592$ ($-32.9\,\%$) & WITH  \\
\midrule
\multicolumn{2}{l}{WITH wins} & \multicolumn{3}{l}{3 of 5 variants (D, SMP, ALL)} \\
\multicolumn{2}{l}{NOSEC wins} & \multicolumn{3}{l}{2 of 5 variants (S\_git, S\_cvss)} \\
\multicolumn{2}{l}{Average $\Delta$} & \multicolumn{3}{l}{$-0.1391$ PR-AUC across the five variants} \\
\bottomrule
\end{tabular}
\caption{GPT ablation, headline ranking metric (test
PR-AUC). The security frontend wins three of five variants
by large margins (26--30 PR-AUC points) and loses two by
smaller margins (3--9 points); the average drop without
the security frontend is 13.9 PR-AUC points across the
five variants.}
\label{tab:gpt-headline}
\end{table}

\begin{table}[!htbp]
\centering\small
\setlength{\tabcolsep}{4.5pt}
\begin{tabular}{l l rr r r}
\toprule
Variant & Metric & WITH & NOSEC & $\Delta$ & Winner \\
\midrule
\emph{D} & PR-AUC    & \textbf{0.8394} & 0.5416 & $-0.2978$ & WITH  \\
\emph{D} & ROC-AUC   & \textbf{0.9400} & 0.8692 & $-0.0708$ & WITH  \\
\emph{D} & F1        & \textbf{0.8095} & 0.5907 & $-0.2188$ & WITH  \\
\emph{D} & Precision & \textbf{0.8293} & 0.5050 & $-0.3243$ & WITH  \\
\emph{D} & Recall    & \textbf{0.7907} & 0.7116 & $-0.0791$ & WITH  \\
\emph{D} & Brier $\downarrow$ & \textbf{0.0493} & 0.1211 & $+0.0718$ & WITH  \\
\midrule
\emph{SMP} & PR-AUC    & \textbf{0.8321} & 0.5738 & $-0.2583$ & WITH  \\
\emph{SMP} & ROC-AUC   & \textbf{0.9549} & 0.9000 & $-0.0549$ & WITH  \\
\emph{SMP} & F1        & \textbf{0.7590} & 0.2257 & $-0.5333$ & WITH  \\
\emph{SMP} & Precision & 0.6655          & \textbf{0.6744} & $+0.0089$ & NOSEC \\
\emph{SMP} & Recall    & \textbf{0.8832} & 0.1355 & $-0.7477$ & WITH  \\
\emph{SMP} & Brier $\downarrow$ & \textbf{0.0657} & 0.1019 & $+0.0362$ & WITH  \\
\midrule
\emph{S\_git} & PR-AUC    & 0.7553          & \textbf{0.7846} & $+0.0293$ & NOSEC \\
\emph{S\_git} & ROC-AUC   & \textbf{0.9250} & 0.9197 & $-0.0053$ & WITH  \\
\emph{S\_git} & F1        & \textbf{0.8065} & 0.5579 & $-0.2486$ & WITH  \\
\emph{S\_git} & Precision & \textbf{0.7576} & 0.4015 & $-0.3561$ & WITH  \\
\emph{S\_git} & Recall    & 0.8621          & \textbf{0.9138} & $+0.0517$ & NOSEC \\
\emph{S\_git} & Brier $\downarrow$ & \textbf{0.0600} & 0.1888 & $+0.1288$ & WITH  \\
\midrule
\emph{S\_cvss} & PR-AUC    & 0.4916 & \textbf{0.5819} & $+0.0903$ & NOSEC \\
\emph{S\_cvss} & ROC-AUC   & 0.8384 & \textbf{0.8718} & $+0.0334$ & NOSEC \\
\emph{S\_cvss} & F1        & 0.4541 & \textbf{0.5701} & $+0.1160$ & NOSEC \\
\emph{S\_cvss} & Precision & 0.5419 & \textbf{0.5636} & $+0.0217$ & NOSEC \\
\emph{S\_cvss} & Recall    & 0.3907 & \textbf{0.5767} & $+0.1860$ & NOSEC \\
\emph{S\_cvss} & Brier $\downarrow$ & 0.1051 & \textbf{0.0975} & $-0.0076$ & NOSEC \\
\midrule
\emph{ALL} & PR-AUC    & \textbf{0.7878} & 0.5286 & $-0.2592$ & WITH  \\
\emph{ALL} & ROC-AUC   & \textbf{0.9390} & 0.8624 & $-0.0766$ & WITH  \\
\emph{ALL} & F1        & \textbf{0.7621} & 0.5568 & $-0.2053$ & WITH  \\
\emph{ALL} & Precision & \textbf{0.6962} & 0.5342 & $-0.1620$ & WITH  \\
\emph{ALL} & Recall    & \textbf{0.8419} & 0.5814 & $-0.2605$ & WITH  \\
\emph{ALL} & Brier $\downarrow$ & \textbf{0.0661} & 0.1013 & $+0.0352$ & WITH  \\
\bottomrule
\end{tabular}
\caption{Full six-metric WITH versus NOSEC comparison for
the five GPT variants on the new 15-run social-media
baseline. Bold marks the better value. Brier is
lower-is-better; all other metrics are higher-is-better.
Per-metric WITH-win tallies (out of five variants): PR-AUC
3, ROC-AUC 4, F1 4, Precision 3, Recall 3, Brier 4. The
single all-NOSEC sweep is \emph{S\_cvss}; the single
all-WITH sweep is \emph{ALL}.}
\label{tab:gpt-full}
\end{table}

\paragraph{What the numbers say.}
The security frontend is doing serious work on the
variants that consume canonical CVE-description language
or the raw social-media post text. On the
description-only variant, removing the security frontend
cuts test PR-AUC by 29.8 points (0.8394 $\to$ 0.5416);
on the combined-everything variant, the drop is 25.9
points (0.7878 $\to$ 0.5286). The new \emph{SMP} variant
shows the most dramatic NOSEC failure mode: PR-AUC drops
by 25.8 points and recall collapses from 0.8832 to
0.1355, meaning the model trained on raw social-media
text without security signals essentially stops firing
positives. These are large effects for a fixed-architecture
ablation and are the clearest evidence yet that the
security overlay is what lets the GNN turn unstructured
text into a reliable exploit-prediction signal.

\paragraph{Where the security frontend does not help, and
            why that makes sense.}
The two NOSEC wins occur on text variants that already
encode the security information in their own way and that
have sparse or redundant security overlays:

\begin{itemize}[leftmargin=*, itemsep=0pt]
    \item \emph{S\_git} is built from a GitHub-URL summary
    that paraphrases commit messages. Commit text rarely
    uses the formal CVE phrasing the regex and dictionary
    layer relies on (``buffer overflow in OpenSSL
    1.0.1f''); it talks about code, diffs, and bug fixes.
    The security frontend contributes a sparse overlay
    here (only 14.9 mean SEC\_* edges per graph, against
    the description's 4.6 baseline and \emph{ALL}'s 57.7;
    see Table~\ref{tab:smp-graph-grid}), and the NOSEC
    model edges past it on PR-AUC and recall while still
    losing badly on precision (0.40 vs 0.76) and F1.
    \item \emph{S\_cvss} is built from a CVSS-metrics
    summary, which already encodes attack vector, attack
    complexity, privileges, scope, and impacts in a fixed
    CVSS vector format. The security frontend's
    overlapping tags become redundant duplicates of
    information the contextual embedding already
    captures: \emph{S\_cvss} produces the lightest
    security overlay of all five variants (only 0.77 mean
    SEC\_* edges per graph; see
    Table~\ref{tab:smp-graph-grid}). Stripping that thin
    overlay actually \emph{helps} the model here: NOSEC
    wins on every one of the six metrics for this
    variant.
\end{itemize}

\paragraph{The calibration story is nearly unanimous.}
Looking only at PR-AUC understates the case. The Brier
score (lower is better) is worse without the security
frontend on four of five variants, and the worst case is
the GitHub-URL summary: Brier balloons from 0.0600 with
security to 0.1888 without, a threefold degradation in
calibration. F1 also drops on four of five variants
without security, with the catastrophic \emph{SMP} F1
collapse (0.7590 $\to$ 0.2257) topping the list. Precision
falls hard on four of five variants (down 32 points on
\emph{D}, down 36 points on \emph{S\_git}, down 16 points
on \emph{ALL}) while recall on those same four either
falls together with it or stays roughly flat, meaning
the NOSEC model is not just wider but truly poorly
calibrated. The security frontend's role is partly to
give the model a strong, narrow signal it can rely on;
strip that signal and the calibration goes with it.

\paragraph{Verdict.}
On the GPT block under the new 15-run baseline the
security frontend wins clearly: the WITH-security model
is better on the three most informative variants
(\emph{D}, \emph{SMP}, \emph{ALL}) by 26--30 PR-AUC
points, and the two losses (\emph{S\_git},
\emph{S\_cvss}) are smaller (3 and 9 PR-AUC points
respectively) and happen on text sources where the
security frontend structurally has little to contribute
(commit-message paraphrase, CVSS-vector restatement).
The average PR-AUC across the five variants drops by
13.9 points when the security frontend is removed --- a
larger effect than the 10.2-point average we reported on
the earlier \emph{S\_all}-based baseline, mostly because
the new \emph{SMP} variant exposes a 26-point WITH-vs-NOSEC
gap that the LLM-curated \emph{S\_all} did not.
Precision, F1, recall, and Brier all degrade in three or
four of the five variants without the security frontend.

The same ablation should still be repeated on the
remaining datasets (the 15 Megavul variants and the
remaining two social-media LLMs) before a final
conclusion is reported. The GPT block is the first slice
on the new baseline and is strongly in favour of keeping
the security frontend; the per-LLM trends summarised in
Section~\ref{sec:per-llm-profile} (Megavul) and
Section~\ref{sec:smp-variant-comparison} (social media)
indicate that Gemma and Mistral are expected to follow
the same pattern.

% ============================================================
\section{Per-LLM characterisation of the Megavul input graphs}
\label{sec:per-llm-profile}

This section answers the natural reviewer question for the
commit-based Megavul corpus: \emph{the security frontend
treats every input the same way, so do the resulting graphs
actually look different across LLM summarisers, and does the
security overlay fire differently?} The social-media half of
this analysis is now consolidated into
Section~\ref{sec:smp-variant-comparison}, which hosts the
15-run social-media baseline end-to-end (graph dimensions,
security overlay, and downstream metrics).

The 15 Megavul (LLM~$\times$~variant) datasets used here
were scanned with two complementary passes:
\begin{itemize}[leftmargin=*, itemsep=0pt]
    \item The processed PyG cache was opened directly and
    the mean number of nodes and edges per graph were
    computed from the per-graph slice offsets. This is
    cheap because no parsing is needed, only tensor reads.
    \item Each labelled record's input text was re-run
    through the rule-only security pipeline, with the
    security-relations layer enabled, to count how often
    each of the ten security entity types and the ten
    typed SEC\_* edges fires.
\end{itemize}

\paragraph{Cross-LLM sanity check.}
On the description-only variant ($D$) every graph is built
from the same CVE description text, so the per-LLM numbers
should agree exactly. They do: every Megavul $D$ run reports
the same 108.6 nodes / 357.6 edges per graph and the same
75.0\,\% of CVEs producing at least one SEC\_* edge. This
confirms the extraction is deterministic and that any
cross-LLM difference observed below is genuinely the LLM
summariser's output, not pipeline noise.

\paragraph{Headline numbers per Megavul LLM.}
\autoref{tab:per-llm-headline} shows the three Megavul blocks
on the most informative variant (\emph{ALL}, which feeds
the model the CVE description plus every available LLM
summary). The columns are: mean nodes per graph, mean edges
per graph, the share of CVEs producing at least one SEC\_*
edge, and the mean number of SEC\_* edges per graph.

\begin{table}[!htbp]
\centering\small
\begin{tabular}{l r r r r r}
\toprule
LLM & Mean nodes & Mean edges & \%CVEs $\geq 1$ SEC & Mean SEC edges & Mean sec.~entities \\
\midrule
GPT      & 848 & 2{,}858 & 99.8\,\% &  57.91 & 22.37 \\
Gemma    & 640 & 2{,}065 & 100.0\,\% & 38.90 & 19.40 \\
Mistral  & 988 & 3{,}201 &  99.9\,\% & 98.83 & 29.06 \\
\bottomrule
\end{tabular}
\caption{Per-LLM characterisation of the Megavul commit-based
graphs on the \emph{ALL} variant. Mean nodes and mean edges
per graph come from a direct read of the processed PyG
cache; the SEC\_* firing rate and the mean counts of typed
security edges and security entities come from re-running
the rule-only security pipeline over each dataset's labelled
records. ``\%CVEs $\geq 1$ SEC'' is the \emph{share} of
test CVEs whose graph contains at least one SEC\_* edge
of any type (it is a coverage measure: a CVE with five
SEC\_* edges and a CVE with one SEC\_* edge both count
once). ``Mean SEC edges'' is the \emph{average count} of
SEC\_* edges per graph (a density measure: that same CVE
with five SEC\_* edges contributes five to the mean).
``Mean sec.\ entities'' sums the ten security entity
categories (CVE\_ID, CWE\_ID, SOFTWARE, VERSION, VULN\_TYPE,
ATTACK\_VECTOR, IMPACT, SEVERITY, REMEDIATION,
CODE\_ELEMENT). The social-media counterpart of this table
has been moved into Section~\ref{sec:smp-variant-comparison},
which now hosts the entire social-media baseline.}
\label{tab:per-llm-headline}
\end{table}

Two patterns stand out on the Megavul block:

\begin{itemize}[leftmargin=*, itemsep=0pt]
    \item \textbf{Mistral generates the largest graphs} (988
    mean nodes, against Gemma's 640). The ratio of largest
    to smallest is around 1.5$\times$.
    \item \textbf{Mistral has by far the most SEC\_* edges
    per graph.} Mistral produces almost twice as many
    SEC\_* edges per graph (98.8) as GPT (57.9), even
    though their mean graph sizes are similar. The
    commit-and-code text Mistral generates is denser in
    canonical-CVE phrasing than the other Megavul LLMs.
\end{itemize}

\paragraph{Per-variant grid.}
\autoref{tab:per-llm-variant-grid} extends the picture to
the 15 Megavul datasets retained after the DeepSeek block
was dropped (its \emph{S\_cvss} variant fell back to the
description text because the DeepSeek source ships with
the CVSS-metrics summary empty for every row, which
breaks per-variant parity across the four LLMs). The
social-media counterpart of this grid is reported in
Section~\ref{sec:smp-variant-comparison}
(Table~\ref{tab:smp-graph-grid}).

One observation that recurs across the grid: the
summary-only \emph{S\_*} variants always produce smaller
and lighter graphs than \emph{ALL}, but they fire the
security overlay almost as often. On GPT's Megavul block,
for example, \emph{S\_url} produces graphs of 175 nodes /
597 edges (against \emph{ALL}'s 848 / 2\,858) but still has
97.2\,\% of CVEs producing a SEC\_* edge against the
\emph{ALL} figure of 99.8\,\%. The LLM summaries on their
own already preserve much of the canonical CVE language
the security frontend keys off.

\begin{table}[!htbp]
\centering\footnotesize
\setlength{\tabcolsep}{6pt}
\begin{tabular}{l l r r r r r}
\toprule
LLM & Variant & \# graphs & Mean nodes & Mean edges & \%CVEs $\geq 1$ SEC & Mean SEC edges \\
\midrule
GPT      & $D$       & 6{,}478 & 109 &   358 & 75.0\,\% &  3.24 \\
GPT      & $S\_url$  & 6{,}474 & 175 &   597 & 97.2\,\% & 14.81 \\
GPT      & $S\_code$ & 6{,}302 & 397 & 1{,}263 & 95.5\,\% & 15.53 \\
GPT      & $S\_cvss$ & 6{,}478 & 195 &   645 & 90.3\,\% & 10.52 \\
GPT      & $ALL$     & 6{,}478 & 848 & 2{,}858 & 99.8\,\% & 57.91 \\
\midrule
Gemma    & $D$       & 6{,}478 & 109 &   358 & 75.0\,\% &  3.24 \\
Gemma    & $S\_url$  & 6{,}478 & 172 &   537 & 97.4\,\% & 11.79 \\
Gemma    & $S\_code$ & 6{,}478 & 233 &   727 & 94.8\,\% & 10.65 \\
Gemma    & $S\_cvss$ & 6{,}478 & 145 &   432 & 89.5\,\% &  7.46 \\
Gemma    & $ALL$     & 6{,}478 & 640 & 2{,}065 & 100.0\,\% & 38.90 \\
\midrule
Mistral  & $D$       & 6{,}478 & 109 &   358 & 75.0\,\% &  3.24 \\
Mistral  & $S\_url$  & 6{,}478 & 208 &   645 & 97.7\,\% & 17.67 \\
Mistral  & $S\_code$ & 6{,}478 & 432 & 1{,}362 & 98.5\,\% & 32.34 \\
Mistral  & $S\_cvss$ & 6{,}362 & 263 &   811 & 87.7\,\% &  9.41 \\
Mistral  & $ALL$     & 6{,}478 & 988 & 3{,}201 & 99.9\,\% & 98.83 \\
\bottomrule
\end{tabular}
\caption{Per-LLM per-variant grid for the 15 retained
Megavul datasets (DeepSeek dropped for missing
CVSS-metrics content). \emph{D} = description only,
\emph{S\_*} = summary-only, \emph{ALL} = description
plus combined summaries. Mean nodes and mean edges come
from the cached graph tensors; the SEC\_* statistics
come from re-running the rule-only security pipeline
over each labelled record. ``\%CVEs $\geq 1$ SEC'' is
the \emph{share} of CVEs whose
graph contains at least one SEC\_* edge of any type
(coverage: 0--100\,\%, each CVE counted once). ``Mean SEC
edges'' is the \emph{average count} of SEC\_* edges per
graph (density: typically 3--99 here, a CVE with many
SEC\_* edges drives the mean up). The two columns
therefore measure different things --- a variant can have
high coverage but low density (e.g.\ \emph{S\_url} fires
some SEC\_* edge on 97.2\,\% of GPT graphs but averages
only 14.8 per graph) or moderate coverage but very high
density (e.g.\ Mistral \emph{ALL} fires on 99.9\,\% of
graphs and averages 98.8 per graph). On every Megavul
row, the \emph{D} variant has both the lowest coverage
(75.0\,\%) and the lowest density (3.24 mean SEC edges)
because the NVD description text in Megavul is short and
code-focused, while every LLM-generated summary
explicitly verbalises attack vectors, impacts, and
software/version references that the security frontend
keys off. The social-media counterpart appears as
Table~\ref{tab:smp-graph-grid} in
Section~\ref{sec:smp-variant-comparison}.}
\label{tab:per-llm-variant-grid}
\end{table}

\paragraph{Which security entities the LLMs surface.}
\autoref{tab:per-llm-entities} breaks the
\emph{ALL}-variant Megavul security entities down by
category. The shared backbone is clear (every LLM surfaces
some SOFTWARE and VERSION mentions, since those phrases
are present in the canonical CVE description), but the
per-LLM ordering of the action-relevant categories differs
sharply. Mistral emits the most VULN\_TYPE, IMPACT, and
CODE\_ELEMENT tags; GPT is the heaviest on ATTACK\_VECTOR
and CODE\_ELEMENT.

\begin{table}[!htbp]
\centering\footnotesize
\setlength{\tabcolsep}{4pt}
\resizebox{\linewidth}{!}{%
\begin{tabular}{l r r r r r r r r r r}
\toprule
LLM & CVE\_ID & CWE\_ID & SOFTWARE & VERSION & VULN\_TYPE & ATTACK\_VEC & IMPACT & SEVERITY & REMED & CODE\_ELEM \\
\midrule
GPT      & 0.05 & 0.00 & 1.98 & 1.69 & 1.35 & 7.61 & 3.02 & 1.76 & 1.03 & 3.87 \\
Gemma    & 0.05 & 0.00 & 1.71 & 1.58 & 2.58 & 4.02 & 3.14 & 2.70 & 0.86 & 2.76 \\
Mistral  & 0.18 & 0.01 & 2.84 & 1.77 & 3.63 & 6.95 & 3.90 & 3.00 & 1.58 & 5.21 \\
\bottomrule
\end{tabular}%
}
\caption{Mean count of each security entity category per
graph on the Megavul \emph{ALL} variant. CWE\_ID
is almost never extracted by the regex because the LLM
summaries rarely quote the literal CWE identifier
pattern; the CVE\_ID counts are also low because the
commit-based text uses the commit hash rather than the
formal CVE identifier. ATTACK\_VECTOR, IMPACT, SEVERITY,
and CODE\_ELEMENT are the categories that vary most
between LLMs and that the model's exploitation-prediction
task cares about. The social-media counterpart appears in
Section~\ref{sec:smp-variant-comparison}
(Table~\ref{tab:smp-entities}).}
\label{tab:per-llm-entities}
\end{table}

\paragraph{Which SEC\_* edges fire.}
\autoref{tab:per-llm-sec-edges} reports the mean count of
each typed SEC\_* edge per Megavul graph on the
\emph{ALL} variant. CAUSES dominates (reflecting the
high IMPACT entity count) and LOCATED\_IN is the second
most common edge because the commit summaries mention
specific code constructs more often than free-text
advisories do.

\begin{table}[!htbp]
\centering\footnotesize
\setlength{\tabcolsep}{4pt}
\resizebox{\linewidth}{!}{%
\begin{tabular}{l r r r r r r r r r r}
\toprule
LLM & AFF & HAS\_V & CLASS & LOC\_IN & EXPL & CAUSES & MITIG & USES & THREAT & HAS\_S \\
\midrule
GPT      &  0.12 &  4.05 & 0.00 &  6.62 & 0.37 & 31.27 & 0.05 & 0.19 & 15.15 & 0.08 \\
Gemma    &  0.10 &  3.22 & 0.00 &  8.68 & 0.18 & 19.35 & 0.04 & 0.12 &  7.09 & 0.12 \\
Mistral  &  0.66 &  5.57 & 0.00 & 23.81 & 1.18 & 45.71 & 0.31 & 0.87 & 20.13 & 0.59 \\
\bottomrule
\end{tabular}%
}
\caption{Mean count of each typed SEC\_* edge per graph
on the Megavul \emph{ALL} variant. Abbreviations: AFF =
AFFECTS, HAS\_V = HAS\_VERSION, CLASS = CLASSIFIED\_AS,
LOC\_IN = LOCATED\_IN, EXPL = EXPLOITED\_BY, MITIG =
MITIGATED\_BY, USES = USES\_FUNCTION, THREAT = THREATENS,
HAS\_S = HAS\_SEVERITY. CLASSIFIED\_AS is near zero
because CWE\_IDs are rarely written in the descriptions
or summaries; CAUSES, LOCATED\_IN, and THREATENS together
account for over 75\,\% of the security overlay on
Megavul because the commit text references code
constructs, impact outcomes, and attack vectors more
explicitly than free-text advisories do. The social-media
counterpart appears as Table~\ref{tab:smp-sec-edges} in
Section~\ref{sec:smp-variant-comparison}.}
\label{tab:per-llm-sec-edges}
\end{table}

\paragraph{Why this matters for downstream performance.}
The PR-AUC ranking inside the Megavul block of the
earlier ablation already echoes these per-LLM
differences. Gemma's Megavul graphs carry the smallest
security overlay (mean of 38.90 SEC\_* edges per graph on
\emph{ALL}), so its model has the shortest auditable
chain back to the text and the weakest SEC view to learn
from. Mistral's Megavul graphs are the largest and
densest (988 nodes, 3{,}201 edges, 98.83 SEC\_* edges per
graph), so the security view contributes the most
clearly there. A reasonable next step for the
methodology section of any paper is to acknowledge this
explicitly: the security frontend is not a free win that
is uniform across LLM summarisers; it is a feature whose
contribution scales with how much canonical CVE language
the LLM keeps in its output. The social-media corpus
(Section~\ref{sec:smp-variant-comparison}) reproduces
this trend in the same direction but with different
absolute magnitudes, dominated by the new \emph{SMP}
variant.

% ============================================================
\section{Social-media variant comparison (15 retrained runs)}
\label{sec:smp-variant-comparison}

A second pass over the social-media corpus retrained every
LLM on a refreshed five-variant set, replacing the earlier
\emph{S\_all} (LLM summary of all sources) with \emph{SMP},
the raw social-media post text itself (truncated to a
16~KB budget; see Section~\ref{sec:inventory}). This
addresses a reviewer concern: the LLM-generated summaries
might be either compressing real signal away or
hallucinating it in. \emph{SMP} sends the underlying post
straight through the pipeline without any LLM in the loop,
so the comparison says \emph{which text source actually
helps the GNN predict exploitation}.

The retrained matrix covers three LLMs (GPT, Gemma,
Mistral) crossed with five variants:
\begin{itemize}[leftmargin=*, itemsep=0pt]
    \item \emph{D} -- the NVD description only.
    \item \emph{SMP} -- the raw social-media post only (no
    LLM summarisation).
    \item \emph{S\_git} -- the LLM's summary of GitHub
    references found for the CVE.
    \item \emph{S\_cvss} -- the LLM's summary of the CVSS
    metric narrative.
    \item \emph{ALL} -- the NVD description plus the three
    LLM summary columns combined.
\end{itemize}
DeepSeek is omitted from this matrix. Its source ships
with the CVSS-metrics summary empty for every row, which
breaks the \emph{S\_cvss} variant and prevents a clean
five-variant comparison; the other three LLMs are
complete. Every run uses the same training schedule as
Section~\ref{sec:gpt-ablation} (multi-view GGNN, hybrid
graph + tabular head, soft-EPSS regression target, EPSS
feature removed from the tabular branch to prevent
leakage).

% ------------------------------------------------------------
\subsection{Graph characterisation of the 15 retrained
datasets}
\label{sec:smp-graph-characterisation}

Before turning to the downstream metrics it is worth
documenting what the GGNN actually consumes. The 15
labelled corpora were re-scanned with the same rule-only
security pipeline used in
Section~\ref{sec:per-llm-profile}, this time reading the
per-CVE labelled records directly (the per-run graph
caches are deleted after evaluation by the disk-cleanup
hook described in Section~\ref{sec:inventory}).

\paragraph{Cross-LLM sanity check.}
On the description-only variant (\emph{D}) every social-media
graph is built from the same NVD description text, so the
per-LLM numbers must match exactly. They do: every \emph{D}
run reports 110.348 mean nodes and 373.904 mean edges per
graph, 4.604 mean SEC\_* edges, and 4.884 mean security
entities. The \emph{SMP} variant is also byte-identical
across LLMs because it bypasses the LLM and feeds the raw
social-media post directly into the pipeline; all three
\emph{SMP} runs report 763.598 nodes, 3{,}374.523 edges,
246.815 SEC\_* edges, and 21.212 security entities per
graph. This confirms the methodology: any cross-LLM
difference on \emph{S\_git}, \emph{S\_cvss}, and
\emph{ALL} is genuinely the LLM summariser's output, not
pipeline noise.

\paragraph{Headline numbers per (LLM, variant).}
\autoref{tab:smp-graph-grid} reports the per-variant
graph size, the security-overlay density, and the
processed-CVE count for every social-media run. \emph{SMP}
is uniformly the densest variant on every LLM (it is the
same input on each row, after all) and 7$\times$ heavier
on nodes than \emph{D}. \emph{ALL} comes in a clear second
because it merges the three LLM summary columns with the
description; the gap from \emph{ALL} to \emph{S\_cvss} or
\emph{S\_git} is large but not as dramatic as the
\emph{SMP} jump.

\begin{table}[!htbp]
\centering\footnotesize
\setlength{\tabcolsep}{5pt}
\begin{tabular}{l l r r r r r}
\toprule
LLM & Variant & \# graphs & Mean nodes & Mean edges & Mean SEC entities & Mean SEC edges \\
\midrule
GPT     & $D$       & 9{,}218 &  110.3 &   373.9 &  4.88 &   4.60 \\
GPT     & $SMP$     & 9{,}212 &  763.6 & 3{,}374.5 & 21.21 & 246.82 \\
GPT     & $S\_git$  & 2{,}246 &  465.4 & 1{,}515.0 & 10.13 &  14.90 \\
GPT     & $S\_cvss$ & 9{,}218 &  188.5 &   606.6 &  5.20 &   0.77 \\
GPT     & $ALL$     & 9{,}218 &  595.8 & 2{,}094.5 & 19.32 &  57.70 \\
\midrule
Gemma   & $D$       & 9{,}218 &  110.3 &   373.9 &  4.88 &   4.60 \\
Gemma   & $SMP$     & 9{,}212 &  763.6 & 3{,}374.5 & 21.21 & 246.82 \\
Gemma   & $S\_git$  & 2{,}381 &  186.4 &   573.8 &  5.79 &   4.64 \\
Gemma   & $S\_cvss$ & 9{,}218 &  140.3 &   412.2 &  2.87 &   0.21 \\
Gemma   & $ALL$     & 9{,}218 &  441.1 & 1{,}430.6 & 14.19 &  31.38 \\
\midrule
Mistral & $D$       & 9{,}218 &  110.3 &   373.9 &  4.88 &   4.60 \\
Mistral & $SMP$     & 9{,}212 &  763.6 & 3{,}374.5 & 21.21 & 246.82 \\
Mistral & $S\_git$  & 2{,}388 &  265.2 &   838.4 &  8.45 &  11.76 \\
Mistral & $S\_cvss$ & 9{,}016 &  265.6 &   811.9 &  5.83 &   0.33 \\
Mistral & $ALL$     & 9{,}218 &  652.9 & 2{,}150.6 & 19.85 &  57.53 \\
\bottomrule
\end{tabular}
\caption{Per-LLM per-variant graph characterisation for the
15 retrained social-media datasets. ``\# graphs'' is the
number of labelled CVEs that produced a graph (drops are
caused by empty source columns: \emph{SMP} loses 6 CVEs
to empty social-media post fields; \emph{S\_git} only
applies to CVEs with a non-empty GitHub-references
summary; Mistral \emph{S\_cvss} loses 202 rows). The
\emph{D} and \emph{SMP} rows are byte-identical across LLMs
because they bypass the LLM. Counts come from re-running
the rule-only security pipeline directly on each labelled
record.}
\label{tab:smp-graph-grid}
\end{table}

\paragraph{Which security entities the LLMs surface.}
\autoref{tab:smp-entities} reports the mean count of each
security entity category on the \emph{ALL} variant. The
shared backbone is the same as the Megavul side: every LLM
surfaces SOFTWARE and VERSION mentions. Mistral is the
heaviest on CVE\_ID, SOFTWARE, and SEVERITY; GPT is the
heaviest on ATTACK\_VECTOR; Gemma is the lightest across
the board, which mirrors its smaller and lighter graphs.

\begin{table}[!htbp]
\centering\footnotesize
\setlength{\tabcolsep}{4pt}
\resizebox{\linewidth}{!}{%
\begin{tabular}{l r r r r r r r r r r}
\toprule
LLM & CVE\_ID & CWE\_ID & SOFTWARE & VERSION & VULN\_TYPE & ATTACK\_VEC & IMPACT & SEVERITY & REMED & CODE\_ELEM \\
\midrule
GPT      & 0.62 & 0.02 & 2.82 & 2.20 & 0.78 & 5.90 & 2.00 & 1.61 & 1.46 & 1.91 \\
Gemma    & 0.73 & 0.02 & 2.74 & 1.84 & 1.03 & 2.73 & 1.58 & 1.77 & 0.88 & 0.88 \\
Mistral  & 1.09 & 0.02 & 3.24 & 2.21 & 1.31 & 3.84 & 1.76 & 3.52 & 1.56 & 1.28 \\
\bottomrule
\end{tabular}%
}
\caption{Mean count of each security entity category per
graph on the social-media \emph{ALL} variant. The
CWE\_ID counts are tiny because the LLM summaries
rarely quote the literal CWE identifier string.
ATTACK\_VECTOR, IMPACT,
SEVERITY, and CODE\_ELEMENT are the categories that vary
most between LLMs and that the exploitation-prediction
task cares about.}
\label{tab:smp-entities}
\end{table}

\paragraph{Which SEC\_* edges fire.}
\autoref{tab:smp-sec-edges} reports the mean count of
each typed SEC\_* edge per graph on the \emph{ALL}
variant. The dominant edges on the social-media corpus
are THREATENS, CAUSES, and HAS\_VERSION, in line with the
high ATTACK\_VECTOR and IMPACT entity counts. Mistral
fires roughly twice as many AFFECTS and HAS\_SEVERITY
edges as the other two LLMs because its summaries quote
software and severity labels more verbatim.

\begin{table}[!htbp]
\centering\footnotesize
\setlength{\tabcolsep}{4pt}
\resizebox{\linewidth}{!}{%
\begin{tabular}{l r r r r r r r r r r}
\toprule
LLM & AFF & HAS\_V & CLASS & LOC\_IN & EXPL & CAUSES & MITIG & USES & THREAT & HAS\_S \\
\midrule
GPT      & 2.53 &  9.31 & 0.01 & 2.45 & 4.21 & 16.93 & 1.05 & 1.44 & 18.53 & 1.24 \\
Gemma    & 2.84 &  7.56 & 0.02 & 1.42 & 2.09 &  6.58 & 0.78 & 0.70 &  8.00 & 1.39 \\
Mistral  & 5.33 & 11.53 & 0.02 & 2.70 & 4.53 & 11.37 & 2.00 & 1.61 & 13.63 & 4.81 \\
\bottomrule
\end{tabular}%
}
\caption{Mean count of each typed SEC\_* edge per graph on
the social-media \emph{ALL} variant. Abbreviations as in
Table~\ref{tab:per-llm-sec-edges}. THREATENS dominates
here because the attack-vector phrases that produce it
(remote, network, attacker) recur in social-media chatter
constantly. CLASSIFIED\_AS is near zero because CWE\_IDs
are rarely written verbatim. The numbers for \emph{SMP}
(not shown) are uniformly much larger across all edge
types because the raw 16\,KB post text contains many more
sentences than any LLM summary; \emph{SMP} averages 246.8
SEC\_* edges per graph against \emph{ALL}'s 31--58.}
\label{tab:smp-sec-edges}
\end{table}

\paragraph{How the graphs map to the downstream metrics.}
The graph characterisation explains two things that show
up later in the downstream tables:
\begin{itemize}[leftmargin=*, itemsep=0pt]
    \item \emph{SMP} produces by far the largest and most
    security-dense graphs (246.8 SEC\_* edges per graph,
    against \emph{ALL}'s 31--58 and \emph{D}'s 4.6). This
    is consistent with \emph{SMP} delivering the highest
    ROC-AUC and recall on every LLM: more SEC\_* anchor
    points give the GNN more lift on the long tail of
    positive CVEs. The flip side is that the extra mass is
    noisier (raw post text, no LLM curation), which is why
    precision is lower and Brier is higher on \emph{SMP}.
    \item \emph{S\_cvss} produces the lightest security
    overlay of all (0.2--0.8 mean SEC\_* edges per graph)
    even though its node and edge counts are healthy. The
    CVSS-metric summary strips out almost every
    natural-language cue that the security frontend keys
    off, leaving the model with effectively only the
    syntactic backbone. This explains why its PR-AUC
    collapses to 0.49--0.55 in
    Table~\ref{tab:smp-full}.
\end{itemize}

% ------------------------------------------------------------
\paragraph{Test-set composition (verified from the saved
per-run test results).}
The \emph{D}, \emph{S\_cvss}, and \emph{ALL} variants
evaluate on the full 1\,385-CVE test split (positive
prevalence 15.5~\%). \emph{SMP} evaluates on 1\,383 CVEs
because five CVEs ship with an empty social-media post
field in the source data and are skipped. \emph{S\_git}
only evaluates on the subset of CVEs that have a
non-empty GitHub-references summary (339--360 test
CVEs, around 24--26~\% of the full
split), so its absolute numbers should be compared against
each other rather than against the other variants.
\autoref{tab:smp-sizes} lists every run's test-set size.

\begin{table}[!htbp]
\centering\small
\setlength{\tabcolsep}{6pt}
\begin{tabular}{l l r r r}
\toprule
LLM     & Variant   & Test CVEs & Positives & Prevalence \\
\midrule
GPT     & D         & 1\,385 & 215 & 15.5\,\% \\
GPT     & SMP       & 1\,383 & 214 & 15.5\,\% \\
GPT     & S\_git    &   339 &  58 & 17.1\,\% \\
GPT     & S\_cvss   & 1\,385 & 215 & 15.5\,\% \\
GPT     & ALL       & 1\,385 & 215 & 15.5\,\% \\
\midrule
Gemma   & D         & 1\,385 & 215 & 15.5\,\% \\
Gemma   & SMP       & 1\,383 & 214 & 15.5\,\% \\
Gemma   & S\_git    &   360 &  59 & 16.4\,\% \\
Gemma   & S\_cvss   & 1\,385 & 215 & 15.5\,\% \\
Gemma   & ALL       & 1\,385 & 215 & 15.5\,\% \\
\midrule
Mistral & D         & 1\,385 & 215 & 15.5\,\% \\
Mistral & SMP       & 1\,383 & 214 & 15.5\,\% \\
Mistral & S\_git    &   359 &  59 & 16.4\,\% \\
Mistral & S\_cvss   & 1\,354 & 213 & 15.7\,\% \\
Mistral & ALL       & 1\,385 & 215 & 15.5\,\% \\
\bottomrule
\end{tabular}
\caption{Test-set composition for the 15 retrained
social-media runs. ``Test CVEs'' is the number of CVEs
in the held-out test split for that run; ``Positives''
is the number of those CVEs whose binary
exploit-prediction label is 1 (i.e.\ EPSS score above
the positive threshold used during training);
``Prevalence'' is the positive rate of the test split,
\emph{Positives / Test CVEs}, expressed as a percentage.
It is the base rate that random and majority-class
baselines would achieve on precision, and it sets the
``no-skill'' floor for PR-AUC --- a classifier that
predicts the prior would score $\text{PR-AUC} \approx
\text{Prevalence}$. The fact that prevalence sits at
roughly 15.5\,\% across all 15 runs means the test sets
are class-imbalanced (only about one CVE in six is
exploitable in the test split), which is why the
downstream tables emphasise PR-AUC and Brier over plain
accuracy. Mistral \emph{S\_cvss} loses 31 CVEs because
the Mistral source ships with an empty CVSS-metrics
summary on a handful of rows.}
\label{tab:smp-sizes}
\end{table}

\paragraph{Headline ranking: test PR-AUC.}
\autoref{tab:smp-headline-pr} is the per-LLM
$\times$ per-variant grid of test PR-AUC. Higher is
better; bold marks the best variant within each row.

\begin{table}[!htbp]
\centering\small
\setlength{\tabcolsep}{8pt}
\begin{tabular}{l r r r r r}
\toprule
LLM     & D & SMP & S\_git & S\_cvss & ALL \\
\midrule
GPT     & \textbf{0.8394} & 0.8321          & 0.7553 & 0.4916 & 0.7878          \\
Gemma   & 0.8312          & 0.8429          & 0.6817 & 0.5513 & \textbf{0.8471} \\
Mistral & 0.8406          & \textbf{0.8408} & 0.5992 & 0.5303 & 0.8317          \\
\bottomrule
\end{tabular}
\caption{Test PR-AUC by LLM and variant for the
15-run social-media ablation. Bold = best per row. The
winning variant differs per LLM (D for GPT, ALL for Gemma,
SMP for Mistral) but the spread is tight: the within-row
gap from best to second-best is at most 0.0073 (Mistral)
and 0.0042 (Gemma).}
\label{tab:smp-headline-pr}
\end{table}

\paragraph{Full six-metric grid.}
\autoref{tab:smp-full} expands the comparison to the
remaining metrics. Brier score is lower-is-better; the
other five metrics are higher-is-better. The two pieces
of evidence that change the story versus
\autoref{tab:smp-headline-pr} alone:
\begin{itemize}[leftmargin=*, itemsep=0pt]
    \item ROC-AUC actually crowns \emph{SMP} for every
    single LLM. \emph{D} and \emph{ALL} win on PR-AUC by
    very small margins on two of the three LLMs, but
    \emph{SMP} consistently delivers the highest ROC-AUC
    (0.9499--0.9572) and the highest recall.
    \item Precision and Brier favour the description-based
    variants (\emph{D} and \emph{ALL}) more clearly.
    \emph{D} achieves 0.83--0.84 precision and the lowest
    Brier of 0.049--0.052, against \emph{SMP}'s 0.67--0.73
    precision and 0.055--0.066 Brier. The model trained on
    raw social-media text is wider (catches more positives
    but with more false positives), the model trained on
    descriptions is narrower.
\end{itemize}

\begin{table}[!htbp]
\centering\small
\setlength{\tabcolsep}{4pt}
\begin{tabular}{l l r r r r r r}
\toprule
LLM & Variant & PR-AUC & ROC-AUC & F1 & Precision & Recall & Brier $\downarrow$ \\
\midrule
GPT     & D       & \textbf{0.8394} & 0.9400          & \textbf{0.8095} & 0.8293          & 0.7907          & \textbf{0.0493} \\
GPT     & SMP     & 0.8321          & \textbf{0.9549} & 0.7590          & 0.6655          & \textbf{0.8832} & 0.0657          \\
GPT     & S\_git  & 0.7553          & 0.9250          & 0.8065          & 0.7576          & 0.8621          & 0.0600          \\
GPT     & S\_cvss & 0.4916          & 0.8384          & 0.4541          & 0.5419          & 0.3907          & 0.1051          \\
GPT     & ALL     & 0.7878          & 0.9390          & 0.7621          & 0.6962          & 0.8419          & 0.0661          \\
\midrule
Gemma   & D       & 0.8312          & 0.9397          & \textbf{0.8096} & 0.8400          & 0.7814          & \textbf{0.0508} \\
Gemma   & SMP     & 0.8429          & \textbf{0.9572} & 0.7699          & 0.7311          & \textbf{0.8131} & 0.0551          \\
Gemma   & S\_git  & 0.6817          & 0.8933          & 0.5275          & 0.7500          & 0.4068          & 0.0974          \\
Gemma   & S\_cvss & 0.5513          & 0.8638          & 0.5673          & 0.5871          & 0.5488          & 0.1005          \\
Gemma   & ALL     & \textbf{0.8471} & 0.9476          & 0.6947          & \textbf{0.8732} & 0.5767          & 0.0573          \\
\midrule
Mistral & D       & 0.8406          & 0.9402          & \textbf{0.7873} & 0.8299          & 0.7488          & \textbf{0.0519} \\
Mistral & SMP     & \textbf{0.8408} & \textbf{0.9499} & 0.7603          & 0.7068          & \textbf{0.8224} & 0.0599          \\
Mistral & S\_git  & 0.5992          & 0.9036          & 0.1231          & 0.6667          & 0.0678          & 0.1311          \\
Mistral & S\_cvss & 0.5303          & 0.8497          & 0.5595          & 0.5270          & 0.5962          & 0.1054          \\
Mistral & ALL     & 0.8317          & 0.9466          & 0.7059          & \textbf{0.8302} & 0.6140          & 0.0566          \\
\bottomrule
\end{tabular}
\caption{Full six-metric grid for the 15 retrained
social-media runs. Bold = best within each LLM block. PR-AUC,
ROC-AUC, F1, Precision and Recall are higher-is-better;
Brier is lower-is-better. SMP wins on ROC-AUC and Recall
across all three LLMs; D wins on F1 and Brier across all
three; PR-AUC and Precision are split per LLM.}
\label{tab:smp-full}
\end{table}

\paragraph{What the numbers say about each text source.}
\begin{itemize}[leftmargin=*, itemsep=0pt]
    \item \textbf{SMP earns its keep.} The raw
    social-media post on its own matches or beats the NVD
    description on the ranking metrics (it wins on ROC-AUC
    for every LLM and is within 0.0073 PR-AUC of the best
    on two of three). This is the headline finding for the
    reviewers who asked whether the post text actually
    carried signal: it does, and it does not need an LLM
    summariser in front of it.
    \item \textbf{The S\_cvss variant is the clear loser
    across the board.} PR-AUC collapses to 0.49--0.55 on
    every LLM, more than 0.28 below the corresponding
    \emph{D} score. The CVSS-metric summary strips the
    natural-language cues that the security frontend keys
    off (and that SecBERT embeds), leaving the GNN with
    little to work with.
    \item \textbf{S\_git lags but on a different test
    set.} \emph{S\_git} evaluates on the 339--360 CVEs
    that actually have GitHub references, with a slightly
    higher positive prevalence (16--17\,\%). Its lower
    PR-AUC is consistent with this smaller, harder slice
    rather than with the GitHub summary being weak in
    absolute terms.
    \item \textbf{ALL is the safest single-variant pick.}
    It is never the worst and is the best for Gemma; for
    GPT and Mistral it is within 0.05 PR-AUC of the
    per-row winner while keeping high precision and a
    small Brier. If the deployment target is one
    variant per LLM rather than per-LLM tuning, ALL is the
    pragmatic default.
\end{itemize}

\paragraph{Verdict.}
The retrained matrix vindicates the new \emph{SMP}
variant and disqualifies \emph{S\_cvss}. The
description-only baseline (\emph{D}) remains a remarkably
strong single-text-source option, especially for
precision-sensitive scoring; \emph{ALL} is the safest
combined baseline. No single text source dominates on
every metric; the trade-off is consistently
\emph{recall-and-ROC} (SMP) versus \emph{precision-and-Brier}
(D / ALL). For the paper, this is exactly the kind of
clean ablation table the reviewers were asking for.

\end{document}
